\section{Rational isometries with the required local theta normalization}
\label{ts-section}

We now produce the isometry required in
Theorem~\ref{tc-kernel-ratio}, with both its local lattice
condition and its absence of fixed vectors.  This removes
the corresponding existence hypothesis from the ordinary
moment criterion.  We give the four-square construction as
well, so that no existence of a rational standard form is
being presumed.

Let $k=\mathbb Q(\sqrt{-D})$, $\mathcal O_k$ and $I$ be
as in Sections~\ref{tc-section} and~\ref{ci-section}.
Write $\sigma=\sqrt{-D}$ for the element with positive
imaginary part in the fixed complex embedding.
Let $W$ be a $k$-vector space of dimension $r$, with
positive definite Hermitian form $h$, linear in its first
variable.  Its associated alternating rational form is
\[
       E_h(x,y)=
           \operatorname{Tr}_{k/\mathbb Q}\!\left(\frac{h(x,y)}{\sigma}\right)
          =\frac2{\sqrt D}\operatorname{Im}h(x,y).
 \tag{TS1}
\]
Let $M\subset W$ be a full $\mathcal O_k$-lattice,
self-dual for $E_h$.

\subsection{The four-square isometry}

\begin{lemma}\label{ts-rational-four-squares}
Every positive rational number is a sum of four rational
squares.
\end{lemma}
\begin{proof}
We include the elementary integer descent.
For an odd prime $\ell$, the pigeonhole argument in
Lemma~\ref{tc-four-squares-congruence} gives integers
$x,y$, chosen with $|x|,|y|\le(\ell-1)/2$, such that
$x^2+y^2+1\equiv0\pmod\ell$.
Thus some $m\ell$ is a sum of four squares with
$0<m<\ell$.
Choose the smallest positive $m$ for which
$m\ell=\sum_{i=1}^4a_i^2$.
If $m>1$, choose integers $b_i\equiv a_i\pmod m$
with $|b_i|\le m/2$.  Then
$\sum_i b_i^2=mm'$ with $0\le m'\le m$.
If $m'=0$, all $a_i$ are divisible by $m$, so
$m\mid\ell$, contradicting $1<m<\ell$.
If $m'=m$, equality in the square bound makes $m$
even and every $a_i=(m/2)(2c_i+1)$.
The sum of four odd squares is divisible by four,
so $m^2\mid\sum_i a_i^2=m\ell$, the same contradiction.
Hence $0<m'<m$.

Regard $a=(a_1,a_2,a_3,a_4)$ and $b=(b_1,b_2,b_3,b_4)$
as integer quaternions.
The four coordinates of $a\bar b$ are divisible by $m$:
modulo $m$, this product is $b\bar b=\sum_i b_i^2=0$.
The multiplicative norm identity, which is the matrix
identity (TC18), now writes
\[
                    m'\ell
                      =\left|\,a\bar b/m\,\right|^2
\]
as a sum of four integer squares, contradicting
minimality of $m$.  Therefore $m=1$.
The prime $2=1^2+1^2$ is immediate.
Multiplying quaternion norms proves the assertion for
every positive integer, by prime factorization.
Finally if $q=A/B>0$ with positive integers $A,B$,
write $AB$ as four integer squares and divide the
four coordinates by $B$.
\end{proof}

\begin{proposition}\label{ts-global-isometry}
There is a $k$-linear isometry
\[
        \gamma_0:(W^{\oplus4},h^{\oplus4})
                           \longrightarrow (k^{4r},h_{\rm st}),
 \qquad h_{\rm st}(x,y)=\sum_jx_j\overline y_j.
 \tag{TS2}
\]
\end{proposition}
\begin{proof}
Hermitian Gram--Schmidt over $k$ diagonalizes $h$ with
positive diagonal entries $q_j\in k\cap\mathbb R=\mathbb Q$.
At each step the orthogonal projection coefficients lie in
$k$, and positivity makes every pivot nonzero.
By Lemma~\ref{ts-rational-four-squares}, write
$q_j=a_j^2+b_j^2+c_j^2+d_j^2$ with rational coordinates.
The quaternion matrix (TC18) for these four coordinates
satisfies $T_j^*T_j=q_jI_4$, since its entries are rational.
After reordering the four copies of the diagonal basis,
the block diagonal matrix with blocks $T_j$ gives (TS2).
\end{proof}

\subsection{Local standard form at a split prime}

\begin{lemma}\label{ts-local-isometry}
If $p$ splits in $k$, there is a $k\otimes\mathbb Q_p$-linear
isometry $\gamma_p$ from $W\otimes\mathbb Q_p$ to the
standard Hermitian space carrying
\[
                     M\otimes\mathbb Z_p
                      \quad\hbox{onto}\quad
                     \mathcal O_k^r\otimes\mathbb Z_p.
 \tag{TS3}
\]
\end{lemma}
\begin{proof}
The split algebra is $\mathbb Q_p\times\mathbb Q_p$,
with conjugation interchanging the factors, and
$\mathcal O_k\otimes\mathbb Z_p=\mathbb Z_p\times\mathbb Z_p$.
Its two idempotents decompose $M\otimes\mathbb Z_p$
as $M_+\oplus M_-$, each free of rank $r$ over
$\mathbb Z_p$.
Hermitian sesquilinearity makes $h$ zero on each one
of the two summands paired with itself.
Write $\sigma=(s,-s)$; the number $s$ is a $p$-adic unit,
since a split prime is unramified.
For $x\in M_+$ and $y\in M_-$, if $h(x,y)=(u,0)$,
then $E_h(x,y)=u/s$.
The perfect alternating pairing on $M\otimes\mathbb Z_p$
therefore says that $h$ gives a perfect integral pairing
between $M_+$ and $M_-$.
Choose a basis of the former and its $h$-dual basis in
the latter.  Sending these to the corresponding standard
bases gives (TS3).  The conjugate cross-pairing is then
correct by the Hermitian identity.
\end{proof}

\begin{theorem}\label{ts-approximation}
Let $\mathcal S$ be a nonempty finite set of split primes
of $k$.  There is a $k$-linear isometry
$\gamma:W^{\oplus4}\to k^{4r}$ such that
\[
 \gamma(M^{\oplus4}\otimes\mathbb Z_p)
          =\mathcal O_k^{4r}\otimes\mathbb Z_p
       \quad(p\in\mathcal S),\qquad
                   \gamma-\gamma_0\ \hbox{is invertible}.
 \tag{TS4}
\]
\end{theorem}
\begin{proof}
Put $m=4r$ and use Lemma~\ref{ts-local-isometry} on
$M^{\oplus4}$ to choose each local $\gamma_p$.
Then $\delta_p=\gamma_p\gamma_0^{-1}$ is unitary on
the standard space.
At a split prime its first component can be any
$G_p\in\operatorname{GL}_m(\mathbb Q_p)$ and its other
component is $(G_p^{\mathsf T})^{-1}$.
The condition that $\delta_p\gamma_0$ carry the specified
lattice onto the standard one is open.
We may perturb $G_p$ within this open set so that
$\det(G_p-1)\det(G_p+1)\ne0$:
replace $G_p$ by $G_p+tI$ for arbitrarily small $t$,
avoiding the finitely many roots of the two determinant
polynomials.

Its Cayley inverse
\[
                     X_p=(\delta_p-1)(\delta_p+1)^{-1}
 \tag{TS5}
\]
is skew-Hermitian and invertible, and $1-X_p$ is
invertible.
The skew-Hermitian matrices over $k$ form a finite-dimensional
$\mathbb Q$-vector space, defined by the rational linear
equation $X^*=-X$.  A rational basis remains a basis after
extension to each $\mathbb Q_p$.
The Chinese remainder theorem, after clearing a finite
set of denominators, approximates the coefficients of
all the $X_p$ simultaneously by coefficients in $\mathbb Q$.
Choose a resulting matrix $X$ sufficiently close that
the determinants of $X$ and $1-X$ stay nonzero and that
\[
                       \delta=(1+X)(1-X)^{-1}
 \tag{TS6}
\]
has the required local lattice property.
To justify the approximation assertion directly, multiply
the finitely many prescribed rational $p$-adic coordinates
by an integer making them integral at the selected primes,
choose their required residues modulo sufficiently high
powers, solve the simultaneous integer congruences, and
divide by that integer.

Since $X^*=-X$, multiplication in (TS6) gives
$\delta^*\delta=1$.
Also $\delta-1=2X(1-X)^{-1}$ is invertible.
The map $\gamma=\delta\gamma_0$ consequently satisfies
both assertions in (TS4).
The zero-dimensional case is immediate with the
usual empty-determinant convention.
\end{proof}

\subsection{Ideal twists and the same determinant correction}

The trace lattices in the source are twisted by the
elliptic ideal.  This twist preserves the hypotheses
above and the local determinant comparison at its
allowed primes.

\begin{lemma}\label{ts-ideal-twist}
The lattice $MI$ is self-dual for $E_h/a$.
If $p\nmid a$, then
\[
                 MI\otimes\mathbb Z_p=M\otimes\mathbb Z_p,
             \qquad I\otimes\mathbb Z_p
                         =\mathcal O_k\otimes\mathbb Z_p.
 \tag{TS7}
\]
Every isometry in (TS2) or (TS4) remains an isometry
from $(h/a)^{\oplus4}$ to $h_{\rm st}/a$.
\end{lemma}
\begin{proof}
The inverse different of $\mathcal O_k$ is
$\sigma^{-1}\mathcal O_k$.
Thus an $\mathcal O_k$-stable lattice self-dual for
the trace form (TS1) is exactly a self-dual Hermitian
lattice: testing $E_h(x,\alpha M)$ for every
$\alpha\in\mathcal O_k$ is equivalent to
$h(x,M)\subset\mathcal O_k$.
This is also seen directly from the trace matrix in
the basis $1,\tau$, whose discriminant is $-D$.
For an invertible ideal,
$(MI)^\#=M^\#\bar I^{-1}$ with respect to $h$,
as follows locally by writing $I$ as a principal ideal.
In a quadratic field $I\bar I=(a)$.
The dual for $h/a$ is consequently
$aM^\#\bar I^{-1}=MI$.
The equivalence with alternating duality proves the
first assertion.
The finite quotient $\mathcal O_k/I$ has order $a$,
which proves (TS7) when $p\nmid a$.
Multiplying both Hermitian forms by $1/a$ does not
change their isometries.
\end{proof}

Let $\kappa$ be the determinant correction of (LM9)--(LM10)
formed from the untwisted lattice $M^{\oplus4}$, reference
isometry $\gamma_0$, and standard target $\mathcal O_k^{4r}$.
For the complex theta construction use the ideal-twisted
lattice $MI$, the Hermitian form $H=(2/(a\sqrt D))h$,
and the standard elliptic lattice $I$.
One may multiply both lattices by the same complex scalar
$\lambda$, replacing each Hermitian form by the transported
form $H_\lambda(\lambda x,\lambda y)=H(x,y)$;
the following conclusion is unchanged.

\begin{theorem}\label{ts-integral-theta}
Fix a split prime $p\ge5$ with $p\nmid a$.
Let $F_M$ be a theta function normalized by
\[
                         F_M^{\otimes4}=\kappa K_{\gamma_0}.
 \tag{TS8}
\]
Such a nonzero theta function exists and is unique up to a
fourth root of unity.
For the single ordinary period $\omega$ of
Proposition~\ref{ci-common-period}, every translate
$U_tF_M$ of order prime to $p$ modulo $MI$
has integral period-scaled moments on $MI$.
The statement also holds on the simultaneously scaled
lattices $\lambda MI$ and $\lambda I$ with the same
period in the lattice directions.
\end{theorem}
\begin{proof}
Apply Theorem~\ref{ts-approximation} with $\mathcal S=\{p\}$.
By (TS7), its isometry $\gamma$ is locally standard for
the ideal-twisted lattices too:
$\gamma((MI)^{\oplus4}\otimes\mathbb Z_p)
=I^{4r}\otimes\mathbb Z_p$.
Their intersection therefore has index prime to $p$.
For finite-index lattices this follows by tensoring the
finite quotient with $\mathbb Z_p$: it vanishes exactly
when the quotient's order is prime to $p$.
The determinant comparison in Theorem~\ref{lm-kappa}
uses the untwisted local equality in (TS4), and gives
\[
                 |\iota(\kappa)|_p^2
                 |\iota(\det(\gamma\gamma_0^{-1}))|_p=1.
 \tag{TS9}
\]
The analytic theorem applies to the ideal-twisted lattices
by Lemma~\ref{ts-ideal-twist}.  Their common norm
semicharacter is exactly the elliptic one by
Lemma~\ref{tc-ideal-character}.
The invertibility in (TS4) gives
$K_{\gamma_0}=cK_\gamma$ with
$c^8=\det(\gamma\gamma_0^{-1})^4$.
Equations (TS9) and this eighth-power equality give
$|\iota(\kappa c)|_p=1$.
Existence and uniqueness in (TS8) follow from the
theta-space argument in Corollary~\ref{tc-normalized-unit}.

By Theorem~\ref{ci-integral-translates}, the product
elliptic theta function and all its prime-to-$p$
translates have integral moments on $I^{4r}$.
The finite intertwiner of index prime to $p$ preserves
this property by Proposition~\ref{lm-finite-intertwiner}.
Pullback by $\gamma$ identifies the direction lattices
exactly.  Multiplying by the unit $\iota(\kappa c)$
therefore gives integral moments for every fourth
tensor power $(U_tF_M)^{\otimes4}$.
The tensor-root theorem, Theorem~\ref{ot-tensor-root},
gives the assertion for $U_tF_M$ itself.
For the scaled lattices, (CI13)--(CI14) identify the
elliptic expansions without any assumption on the
absolute value of $\iota(\lambda)$; the same finite
operator and pullback proof then applies.
\end{proof}

This theorem supplies the ordinary theta moments for the
normalized self-dual lattice vector itself, with the same
determinant correction as the untwisted trace lattice and
for every split prime $p\geq5$ prime to the elliptic
ideal norm.  This supplies all sufficiently large
selected split primes.  The source's generic moment
statement also permits $p=3$; that case is outside
the Bannai--Kobayashi input used in this reconstruction.
It does not yet identify the more general arithmetic
column, which includes further finite lattice moves,
local support conditions and incidence maps.


